%% =====================================================================
%%  B7-T-23-07 -- what B7 contributes to "The Wrath of the Liar"
%%  -------------------------------------------------------------------
%%  LAYER A: APPEND-ONLY. Written once at the entry's write, then left
%%  alone. If the source has more to say later, add a bullet -- do not
%%  rewrite. Material found ONLY here goes in the `unique` slot with
%%  the unique flag set, which makes it undeletable (rule R10).
%% =====================================================================
%% @kind: source
%% @id: B7-T-23-07
%% @book: B7
%% @topic: T-23-07
%% @locator: ch. 7 (pp. 51--56), with ch. 5 attack-mode items (pp. 40--44)
%% @status: distilled
%% @unique: no
%% @integrated: yes
%% @updated: 2026-09-19

\dossier{B7}{T-23-07}{\bookshort{B7}}{ch. 7 (pp. 51--56), with ch. 5 attack-mode items (pp. 40--44)}

\begin{distilled}
  \item When facts corner a subject, attack mode is the most available verbal indicator: the counterattack on the questioner or the process substitutes for an answer (the Enron CEO's 'witch hunt' testimony as the extended case).
  \item The attack family: going into attack mode; process or procedural complaints (objection to how replaces answer to what); inappropriate levels of politeness or concern toward the questioner; sudden failure to understand a simple question.
  \item Nietzsche's epigraph on the chapter: no man lies so boldly as the man who is indignant.
  \item B7 treats the display as data within the timing model --- an indignation flare inside the five-second window counts toward the cluster; outside it, or against a hot-tempered baseline, it is only temperament.
  \item The pressure function: the counterattack puts the questioner on the defensive about tone, and most people abandon the original question there --- which trains the operator to lead with outrage.
\end{distilled}
